
We decomposed the problem of estimating the ski jumper’s trajectory into a sequence of well-defined sub-tasks,
so that each component could be designed, tested, and improved independently.

\medskip

First, we performed per-frame detection of the athlete using YOLO.
From the detected bounding box, we extracted a region of interest that isolates the jumper from the background
and reduces the influence of irrelevant image content.

\medskip

Next, we applied SAM3 to obtain a pixel-level segmentation of the jumper.
From this segmentation mask, we estimated the center of mass (CoM) of the athlete,
which we used as the main measurement to represent the jumper position over time
and to build an initial (raw) trajectory.

\medskip

In early experiments, the CoM estimates exhibited temporal jitter and occasional outliers,
mainly due to segmentation noise, partial occlusions, and motion blur.
To enforce temporal coherence and reduce instability, we introduced a Kalman filter.
The filtered CoM sequence provides a smoother and more reliable observation stream,
while still preserving the underlying dynamics of the motion.

\medskip

Since the input videos are multi-camera and include abrupt viewpoint changes,
we explicitly detected camera switches to avoid mixing incompatible geometries.
We used an HSV-based algorithm to identify frames where the color distribution changes sharply,
which is a strong indicator of a cut.
Using these detections, we split the full video into homogeneous segments (shots)
in which the viewpoint is approximately consistent.

\medskip

For each shot, we selected as reference the first frame in which the jumper is visible.
Starting from this reference, we performed synchronization-based image mosaicking
to align frames and express all measurements in a common coordinate system.
This step is crucial to compensate for camera motion and to compare the jumper position across time.

\medskip

We formulated the alignment problem as a graph optimization task.
Each frame in the shot is represented as a node,
and each edge encodes a planar affine transformations that maps one frame into another.
At first, we experimented with a homography-based approach, 
but it proved too unstable due to the lack of sufficient keypoints in each scene and the presence of significant motion blur. 
Therefore, we decided to adopt an affine transformation instead.
To increase robustness and mitigate drift, we did not rely only on consecutive correspondences:
for each frame $x$, we estimated affine transformation both between $x$ and $x-1$
and between $x$ and $x-4$.
The longer-baseline edges help stabilize the solution when consecutive frames provide weak matches,
for instance due to motion blur or limited texture.

\medskip

All estimated affine transformation were constrained through an explicit consistency condition,
so that transformations remain compatible when composed along different paths in the graph.
In practice, this enforces that multiple routes connecting the same pair of frames
produce coherent mappings, reducing accumulated error and improving global alignment.

\medskip

The resulting system exhibits gauge freedom,
meaning that without an additional constraint it admits multiple equivalent solutions
(e.g., any global change of reference can produce the same relative affine transformation).
To remove this ambiguity, we fixed the gauge by introducing a gauge weight
that anchors the solution to a chosen reference and ensures uniqueness.

\medskip

Finally, we solved the resulting overdetermined system using least squares,
obtaining the minimum-norm solution that best satisfies all pairwise homography constraints.
Once the frames are globally aligned, the (filtered) CoM measurements can be consistently projected
into the reference coordinate system, yielding the final trajectory for each shot.